% Phase 15 only: complete comparative statics and mandatory boundary checks.
\section{Comparative Statics and Boundary Cases}
\label{sec:p15-comparative-statics}

This section distinguishes four operations that cannot be interchanged:
\begin{enumerate}
  \item a derivative with respect to a continuous primitive at a fixed
  candidate CMO price \(p_m\);
  \item a derivative with respect to a continuous primitive after the CMO
  market clears at \(p_m^*\);
  \item a finite comparison between the binary regimes \(M=0\) and \(M=1\);
  \item a limiting-case calculation.
\end{enumerate}
There is no derivative with respect to binary \(M\). Before each economic
interpretation, the mathematical margin is named as commercial return,
commercialization-route choice, project advancement, the CMO price, or
realized product output. Upstream research and patent generation are not
endogenous margins of the baseline.

\subsection{Commercial pricing and return}
\label{subsec:p15-commercial-return}

\paragraph{Object and held-fixed variables.}
The differentiated objects are the optimized product price, operating profit,
and commercial present value. Project value \(q\), marginal manufacturing cost
\(c\), and the demand and discount primitives are held fixed except for the
single argument named in each derivative. These are conditional product-market
calculations; neither \(M\) nor \(p_m^*\) is varied.

The Phase~2 solution and operating profit can be written as
\[
  p^*(c)=\frac{\varepsilon}{\varepsilon-1}c,
  \qquad
  \pi(q,c)
  =
  \frac{Aq}{\varepsilon}
  \left(\frac{\varepsilon}{\varepsilon-1}c\right)^{1-\varepsilon}.
\]
Direct differentiation gives
\begin{equation}
  \frac{\partial p^*}{\partial c}
  =
  \frac{\varepsilon}{\varepsilon-1}>1,
  \qquad
  \frac{\partial\pi}{\partial q}
  =
  \frac{\pi}{q}>0,
  \qquad
  \frac{\partial\pi}{\partial c}
  =
  (1-\varepsilon)\frac{\pi}{c}<0.
  \label{eq:p15-price-profit-derivatives}
\end{equation}
Since \(R(q,c)=\pi(q,c)/(1-\beta\varphi)\),
\begin{equation}
  R_q=\frac{R}{q}>0,
  \qquad
  R_c=(1-\varepsilon)\frac{R}{c}<0,
  \label{eq:p15-return-derivatives}
\end{equation}
and, holding \(\pi\) fixed,
\[
  R_\beta
  =
  \frac{\varphi\pi}{(1-\beta\varphi)^2}\geq0,
  \qquad
  R_\varphi
  =
  \frac{\beta\pi}{(1-\beta\varphi)^2}>0.
\]
The first inequality is an equality when \(\varphi=0\).

\paragraph{Sufficient conditions, dimensions, and zero effects.}
The signs use \(A,q,c>0\), \(\varepsilon>1\),
\(\beta\in(0,1)\), and \(\varphi\in[0,1)\). The price derivative is
dimensionless; \(R_q\) and \(R_c\) have currency per unit of their respective
arguments. A change in \(\beta\) has no continuation-value effect when
\(\varphi=0\). None of these derivatives is a policy effect unless an approved
route technology changes the cost argument supplied to \(R\).

\paragraph{Economic interpretation and empirical prediction.}
This is a commercial-return margin. Higher manufacturing cost is passed
through to the product price and lowers optimized commercial value, while a
higher project-value shifter raises it. The empirical prediction is conditional
cost pass-through and lower commercial return for more costly manufacturing,
not a claim about patenting or upstream research.

\subsection{Commercialization-route choice at fixed CMO price}
\label{subsec:p15-route-choice}

\paragraph{Object and held-fixed variables.}
The object is the internal-minus-entrusted value gap and its conditional
capability cutoff. Hold \((q,m)\), all technologies and outside options fixed.
For capability derivatives also hold finite \((\tau_E,p_m)\) fixed. For each
implicit cutoff derivative, vary only the continuous argument displayed.

On the internally feasible domain,
\begin{equation}
  \Delta_{IE,k}
  =
  s(q)R_c(q,c_I)c_{I,k}-F_{I,k}>0.
  \label{eq:p15-gap-capability}
\end{equation}
The operating-cost component is weakly positive and may be zero when
\(s(q)=0\); the setup-cost component is strictly positive. Under the
cell-specific crossing condition, the implicit-function theorem gives
\begin{equation}
  \left.\frac{\partial k^*}{\partial\tau_E}\right|_{q,m,p_m}
  =
  -\frac{1}{\Delta_{IE,k}(k^*)}<0,
  \qquad
  \left.\frac{\partial k^*}{\partial p_m}\right|_{q,m,\tau_E}
  =
  -\frac{b(m)}{\Delta_{IE,k}(k^*)}<0.
  \label{eq:p15-cutoff-derivatives}
\end{equation}
Both are continuous-argument derivatives inside the finite-wedge domain.
They are not derivatives with respect to \(M\), and the second is not an
equilibrium-price effect.

\paragraph{Binary policy comparison.}
Let \(W_i^0=\max\{W_i^I,W^T,0\}\). At any fixed finite \(p_m\),
\begin{equation}
  W_i(q,m;1,p_m)-W_i(q,m;0,p_m)
  =
  [W_i^E(q,m;1,p_m)-W_i^0(q,m)]_+\geq0.
  \label{eq:p15-fixed-project-gain}
\end{equation}
This finite comparison adds an option without changing the old route values.
It is strict only on \(\mathcal C_i(p_m)\).

\paragraph{Sufficient conditions and zero effects.}
The cutoff signs require internal feasibility, the Phase~3 technology signs,
finite \((\tau_E,p_m)\), strict \(\Delta_{IE,k}>0\), and endpoint crossing.
If crossing fails, no interior cutoff is asserted. If \(T\) or \(A\) dominates,
the \(I/E\) cutoff does not determine the realized route. If
\(W_i^E\leq W_i^0\), the finite policy gain is zero.

\paragraph{Economic interpretation and empirical prediction.}
This is a commercialization-route margin. Better internal manufacturing
capability makes internal production relatively more attractive; a lower
finite entrusted-route burden expands the capability region using \(E\), while
a higher fixed CMO price contracts it. The empirical prediction is sorting of
holder--manufacturer organization by predetermined manufacturing capability
and project requirement, conditional on outside options.

\subsection{Project advancement at fixed CMO price}
\label{subsec:p15-advancement-fixed-price}

\paragraph{Object and held-fixed variables.}
The object is the unique common project-advancement intensity
\[
  x_i^*(M,p_m)
  =
  \left[\frac{\beta a_i\Omega_i(M,p_m)}{\kappa}\right]^{1/\nu}.
\]
Each partial derivative below holds \(M,p_m,\nu\) and all unlisted primitives
fixed. In particular, route-value objects are held fixed when \(a_i\) or
\(\kappa\) is varied. The formulas apply at \(\Omega_i>0\); the
\(\Omega_i=0\) corner is handled separately.

Differentiation of the Phase~5 FOC gives
\begin{equation}
  \frac{\partial x_i^*}{\partial\Omega_i}
  =
  \frac{x_i^*}{\nu\Omega_i}>0,
  \qquad
  \frac{\partial x_i^*}{\partial a_i}
  =
  \frac{x_i^*}{\nu a_i}>0,
  \qquad
  \frac{\partial x_i^*}{\partial\kappa}
  =
  -\frac{x_i^*}{\nu\kappa}<0.
  \label{eq:p15-advancement-partials}
\end{equation}
The corresponding elasticities are \(1/\nu,1/\nu,-1/\nu\).
Because the optimized project value obeys the deterministic envelope,
\(\Omega_{i,p_m}=-\chi_i^E\) almost everywhere, so
\begin{equation}
  \left.\frac{\partial x_i^*}{\partial p_m}\right|_M
  =
  -\frac{x_i^*}{\nu\Omega_i}\chi_i^E
  \leq0.
  \label{eq:p15-advancement-price}
\end{equation}
This is a fixed-candidate-price derivative and already includes both the
intensive value effect and deterministic route-envelope effect inside
\(\Omega_i\).

For binary \(M\), define at a common support price
\[
  \Delta\Omega_i(p_m)
  =
  \int[W_i^E(q,m;1,p_m)-W_i^0(q,m)]_+\,dF.
\]
The correct policy calculation is the finite difference
\begin{equation}
  \Delta x_i(p_m)
  =
  \left(\frac{\beta a_i}{\kappa}\right)^{1/\nu}
  \left[
    (\Omega_i^0+\Delta\Omega_i(p_m))^{1/\nu}
    -(\Omega_i^0)^{1/\nu}
  \right]
  \geq0.
  \label{eq:p15-fixed-advancement-gain}
\end{equation}
It is strict if and only if \(\Delta\Omega_i(p_m)>0\).

\paragraph{Sufficient conditions and zero effects.}
The partials require \(\Omega_i>0\). At the zero-value corner, \(x_i^*=0\);
global monotonicity follows directly from the closed form. Higher \(a_i\)
scales a positive response but cannot create one when the entrusted surplus is
zero. The weak price effect is zero when no project uses \(E\), so
\(\chi_i^E=0\). The manufacturing-capability ordering of the level response
additionally uses \(\nu\geq1\); no such sign is claimed for
\(0<\nu<1\) without another bound.

\paragraph{Economic interpretation and empirical prediction.}
This is development-stage project advancement, not upstream research.
Expected commercialization value raises the return to advancing viable
projects; higher advancement productivity amplifies a positive value gain,
whereas a steeper cost scale reduces the chosen intensity. The prediction is
a larger project-advancement response among high-\(a_i\) developers for which
entrusted manufacturing adds value, especially when internal manufacturing
capability is weak under the stated curvature condition. No patent-application
response follows mechanically.

\subsection{CMO supply, demand, and continuous equilibrium derivatives}
\label{subsec:p15-cmo-continuous}

\paragraph{Supplier object and held-fixed variables.}
For a supplier with \(p_m>0\), vary only candidate price or efficiency while
holding the other argument and the cost function fixed. Implicit
differentiation of \(p_m=\Psi_s(s_j^*;z_j)\) gives
\begin{equation}
  \frac{\partial s_j^*}{\partial p_m}
  =
  \frac{1}{\Psi_{ss}}>0,
  \qquad
  \frac{\partial s_j^*}{\partial z_j}
  =
  -\frac{\Psi_{sz}}{\Psi_{ss}}>0.
  \label{eq:p15-supplier-partials}
\end{equation}
These are capacity-supply derivatives; policy does not shift
\(\Psi\) or \(z_j\).

\paragraph{Study-demand object.}
At differentiability points of the aggregated deterministic-choice functions,
\begin{equation}
  \frac{\partial D_m^{\mathrm{MAH}}}{\partial p_m}
  =
  \int a_i
  \left[
    \frac{\partial x_i^*}{\partial p_m}\chi_i^E
    +
    x_i^*\frac{\partial\chi_i^E}{\partial p_m}
  \right]dH(a,k)
  \leq0.
  \label{eq:p15-study-demand-slope}
\end{equation}
The first term is the project-value/advancement response; the second is the
commercialization-route response. Where a classical derivative does not
exist, the global weak decrease follows from monotonicity and atomless
aggregation.

\paragraph{Continuous equilibrium perturbation.}
Consider an infinitesimal continuous, non-policy perturbation. Let
\(dD_m\) and \(dS_m\) denote its direct demand and supply shifts at a fixed
price. Hold every other primitive fixed. Differentiating
\(Z(p_m)=D_m(p_m)-S_m(p_m)=0\) at the unique equilibrium gives
\(S_{m,p}-D_{m,p}>0\) and
\begin{equation}
  dp_m^*
  =
  \frac{dD_m-dS_m}
       {S_{m,p}-D_{m,p}}.
  \label{eq:p15-equilibrium-ift}
\end{equation}
A demand-increasing shifter raises the equilibrium price when supply is not
shifted; a supply-increasing shifter lowers it when demand is not shifted.
This formula is not applied to binary \(M\).

\paragraph{Sufficient conditions, dimensions, and zero effects.}
The supplier signs use \(\Psi_{ss}>0\) and \(\Psi_{sz}<0\). The aggregate
price derivative uses differentiability and the same strict excess-demand
slope that ensures uniqueness. It is zero when the perturbation changes
neither demand nor supply at the equilibrium. Equation
\eqref{eq:p15-equilibrium-ift} is dimensionally a change in CMO price.

\paragraph{Interpretation and prediction.}
These are CMO capacity and equilibrium-price margins. A higher capacity price
induces more qualified supply but reduces study demand through both fewer
entrusted assignments and weaker project advancement. The prediction is an
upward-sloping supply response and a scarcity price that responds positively
to continuous demand pressure, conditional on unchanged supply technology.

\subsection{Binary reform at the equilibrium CMO price}
\label{subsec:p15-equilibrium-policy}

\paragraph{Object and finite comparison.}
Let \(p_m^h=p_m^*(h)\) for \(h\in\{0,1\}\). Supply technology, the supplier
distribution, and background demand are held fixed. Since
\(D_m^{\mathrm{MAH}}(p_m;0)=0\), market clearing at \(p_m^0\) implies
\[
  Z_1(p_m^0)=D_m^{\mathrm{MAH}}(p_m^0;1)\geq0.
\]
Strict monotonicity of \(Z_1\) therefore gives the finite comparison
\begin{equation}
  p_m^1\geq p_m^0,
  \qquad
  p_m^1>p_m^0
  \ \Longleftrightarrow\
  D_m^{\mathrm{MAH}}(p_m^0;1)>0.
  \label{eq:p15-equilibrium-price-order}
\end{equation}
No derivative with respect to \(M\) is used.

Define the fixed-price direct gain and the equilibrium gain as
\[
  \Delta\Omega_i^{\mathrm{dir}}
  =
  \int[W_i^E(q,m;1,p_m^0)-W_i^0(q,m)]_+\,dF,
\]
\[
  \Delta\Omega_i^{\mathrm{eq}}
  =
  \int[W_i^E(q,m;1,p_m^1)-W_i^0(q,m)]_+\,dF.
\]
Because \(W_{i,p_m}^E=-b(m)<0\), the price order yields
\begin{equation}
  0\leq\Delta\Omega_i^{\mathrm{eq}}
  \leq\Delta\Omega_i^{\mathrm{dir}}.
  \label{eq:p15-value-attenuation}
\end{equation}
Applying the strictly increasing advancement mapping to these value gains
gives
\begin{equation}
  0\leq\Delta x_i^{\mathrm{eq}}
  \leq\Delta x_i^{\mathrm{dir}}.
  \label{eq:p15-advancement-attenuation}
\end{equation}

\paragraph{Conditions and zero effects.}
The price order requires policy-invariant supply/background demand and
nonnegative post-reform study demand. Strictness requires positive study
demand at \(p_m^0\). A developer's equilibrium gain is zero when the
entrusted option remains below its old-route maximum at \(p_m^1\). Scarcity
can eliminate an incremental gain but cannot lower the old choice-set value.

\paragraph{Economic interpretation and empirical prediction.}
This is an equilibrium CMO-price feedback on project advancement. New
entrusted-route demand bids up scarce qualified capacity, so a fixed-price
calculation overstates the equilibrium value and advancement response. The
prediction is a weak post-reform increase in the CMO capacity price when the
study cohort creates positive demand and a smaller project response in markets
with less elastic capacity supply.

\subsection{Planning-stage and realized-product finite comparisons}
\label{subsec:p15-outcomes}

\paragraph{Objects and held-fixed primitives.}
The objects are planning-stage project intensity and realized retained-product
outcomes, evaluated at the regime-specific equilibrium. Predetermined
\(a_i\), project distributions, realization probabilities, technologies and
outside options are held fixed. The regime changes only route availability and
the endogenous CMO price.

Planning-stage intensity satisfies the exact finite difference
\begin{equation}
  \Delta\Lambda_i^{\mathrm{plan}}
  =
  a_i\Delta x_i^{\mathrm{eq}}\geq0.
  \label{eq:p15-planning-change}
\end{equation}
Let
\[
  Q_i^{\mathrm{ret},h}
  =
  \int s(q)
  [\mathbf 1\{r_i^{*,h}=I\}+\mathbf 1\{r_i^{*,h}=E\}]\,dF.
\]
Adding route \(E\) cannot remove an old option. If \(E\) wins, it is retained;
if it does not win, the old maximizer is unchanged. Hence
\(Q_i^{\mathrm{ret},1}\geq Q_i^{\mathrm{ret},0}\). The exact decomposition is
\begin{equation}
  \Delta Y_i^{\mathrm{ret}}
  =
  a_i(x_i^1-x_i^0)Q_i^{\mathrm{ret},0}
  +
  a_ix_i^1
  (Q_i^{\mathrm{ret},1}-Q_i^{\mathrm{ret},0})
  \geq0.
  \label{eq:p15-retained-outcome-change}
\end{equation}
The first term is project advancement; the second is retained-route
composition. The same decomposition applies within class \(g\) using
\(\rho_gF_g\) and the common \(x_i\), not a class-specific control.

\paragraph{Sufficient conditions and zero effects.}
A strict planning-stage increase requires a positive equilibrium expected-value
gain. A strict retained-product increase requires a positive advancement term
or positive-measure movement from \(T/A\) into \(E\) with positive realization
weight. Reassignment from \(I\) to \(E\) raises observed holder--producer
separation but, by itself, does not change the retained total. If advancement
and route indicators are unchanged, every displayed outcome change is zero.

\paragraph{Economic interpretation and empirical prediction.}
These are planning-stage and realized-product margins. The model predicts that
retained holder--producer separation can rise after route \(E\) becomes
available, and separates that assignment change from the number of viable
projects advanced and from downstream realization. Product-level holder and
manufacturer records can test the assignment prediction. Patent applications
are not an outcome in these expressions.

\subsection{Mandatory limiting-case audit}
\label{subsec:p15-boundaries}

\paragraph{Internal capability \(k_i\to\infty\).}
On the feasible domain, \(\Delta_{IE,k}>0\), so increasing capability weakly
favors \(I\) relative to \(E\). Under the maintained upper-end crossing
condition, \(\Delta_{IE}>0\) for sufficiently large \(k_i\). This is a relative
route-value statement; \(T\) or \(A\) may still dominate both.

\paragraph{Prohibitive internal capability.}
If \(k_i<\underline{k}(m)\), then \(F_I(m,k_i)=+\infty\) and
\(W_i^I=-\infty\). Route \(I\) is infeasible rather than merely unattractive.

\paragraph{Pre-MAH regime.}
When \(M=0\), \(\tau_E(0)=+\infty\), so \(W_i^E=-\infty\),
\(\chi_i^E=0\), and \(D_m^{\mathrm{MAH}}=0\). The old route set and background
CMO market remain.

\paragraph{Infinite CMO price.}
Because \(b(m)>0\),
\[
  \lim_{p_m\to\infty}W_i^E(q,m;1,p_m)=-\infty.
\]
Thus \(E\) ceases to be optimal and
\begin{equation}
  \lim_{p_m\to\infty}
  W_i(q,m;1,p_m)
  =
  W_i^0(q,m).
  \label{eq:p15-infinite-price-boundary}
\end{equation}
The newly added option disappears economically but cannot reduce the old
route-set value.

\paragraph{Zero entrusted-route advantage.}
If \(W_i^E(q,m;1,p_m)\leq W_i^0(q,m)\) for every draw, then
\begin{equation}
  \Delta\Omega_i(p_m)=0,
  \qquad
  \Delta x_i(p_m)=0.
  \label{eq:p15-zero-advantage-boundary}
\end{equation}
At the equilibrium price, the same conclusion applies when the inequality
holds at \(p_m^1\).

\paragraph{Perfectly elastic CMO supply.}
If qualified capacity is available at the common administered/competitive
price \(\bar p_m\) in both regimes, then
\[
  p_m^1=p_m^0=\bar p_m,
  \qquad
  \Delta\Omega_i^{\mathrm{eq}}=\Delta\Omega_i^{\mathrm{dir}},
  \qquad
  \Delta x_i^{\mathrm{eq}}=\Delta x_i^{\mathrm{dir}}.
\]
The equilibrium attenuation channel disappears. This is a boundary outside
the strictly increasing-supply condition used for uniqueness, not a claim
that baseline supply is perfectly elastic.

\paragraph{Quadratic advancement cost.}
For \(\nu=1\),
\begin{equation}
  x_i^*(M,p_m)
  =
  \frac{\beta a_i}{\kappa}\Omega_i(M,p_m),
  \qquad
  \Delta x_i
  =
  \frac{\beta a_i}{\kappa}\Delta\Omega_i.
  \label{eq:p15-quadratic-boundary}
\end{equation}
The value response is exactly linear.

\paragraph{Novelty-share boundaries and no fixed ranking.}
The common-control decomposition is
\begin{equation}
  \Delta\Omega_i
  =
  \rho_O\Delta\Omega_{iO}
  +
  \rho_{\mathrm{Inc}}\Delta\Omega_{i\mathrm{Inc}}.
  \label{eq:p15-novelty-boundary}
\end{equation}
If \(\rho_O=0\), the original-oriented class contributes zero to the aggregate;
if \(\rho_{\mathrm{Inc}}=0\), the incremental class contributes zero. Either
conditional gain may be zero while the other is positive. Reversing the class
supports of positive entrusted surplus reverses their ordering, so no fixed
ranking follows from the baseline.

\subsection{Comparative-static scope conclusion}
\label{subsec:p15-scope}

All central signs above follow from the four Phase~14 assumption blocks and
the deterministic option-value structure. The module adds no route, control,
state, market, welfare object, patent outcome, logit share, inclusive value or
continuous implementation index. Its finite policy chain is
\[
  \begin{aligned}
  M
  &\longrightarrow
  \text{availability of retained entrusted manufacturing}
  \longrightarrow
  \Omega_i \\
  &\longrightarrow
  x_i^*
  \longrightarrow
  \text{planning, route, and realized-product objects}.
  \end{aligned}
\]
with the endogenous CMO price attenuating, rather than directly causing, the
project-value response.
